\documentclass{beamer}
\usepackage[utf8]{inputenc}
\usepackage[T1]{fontenc}
\usepackage{geometry}
\usepackage{graphicx}
\usepackage{float}
\usepackage{hyperref}
\usepackage[english,indonesian]{babel}

\title{Optimasi Pemrosesan Template LaTeX Menggunakan Pendekatan Berbasis AI}
\author{Academic Support Agent}
\date{\today}

\begin{document}
\maketitle

\begin{abstract}

\end{abstract}

\chapter{PENDAHULUAN}


\chapter{TINJAUAN PUSTAKA}


% Original content (commented out):
% \documentclass[10pt,aspectratio=169]{beamer}
\usepackage[brazil]{babel}
\usepackage{source/littoral}
\usepackage[none]{hyphenat}

\title{\textbf{Título da Apresentação}}
\subtitle{Subtítulo}
\author{Nome Sobrenome}

\definecolor{verde}{HTML}{23B037}
\definecolor{verm}{HTML}{D02337}

\begin{document}

\maketitle

\section{Introdução}

\begin{frame}{\thesection \, \secname}{Contextualização}
	\begin{itemize}[<+->]
		\item Tópico 1
		\item Tópico 2
	\end{itemize}
\end{frame}

 
\begin{frame}{\thesection \, \secname}{Problemática}
	\begin{itemize}[<+->]
		\item Tópico 1
		\item Tópico 2
	\end{itemize}
\end{frame}

\section{Metodologia}
\begin{frame}{\thesection \, \secname}{Parametrização}
	\begin{itemize}
		\item Tópico 1
		\item Tópico 2
	\end{itemize}
	\begin{figure}[hbtp]
	\centering
	\includegraphics[scale=0.3]{figs/met1.png}
	\end{figure}	
\end{frame}


\section{Conclusão}
\begin{frame}{\thesection \, \secname}{}
	\begin{itemize}
		\item Tópico 1
		\item Tópico 2
	\end{itemize}
\end{frame}




\backmatter

\end{document}


\end{document}